\section{Experimental Setup}\label{sec:setup}

\textbf{Benchmark.} We use the test methods of the UTRefactor replication package~\cite{utrefactor}
(Apache Commons-cli, Gson, Jsoup). The package has reports for 492 test methods; 484 match the project commits we pinned
verbatim and form our evaluation set (Commons-cli 77, Gson 238, Jsoup 169; the pinned commits contain 97.9--98.8\% of the published inputs verbatim).
Projects are built at the pinned commit with Maven~3.9.9 (each project with its pinned JDK~17 where its build needs it; JNose and the tool chain run on JDK~21); Gson is evaluated against its own test dependencies at that commit.

\textbf{Splits.} The 484 tests are split once, stratified by project and single/multi-smell, with a fixed seed (commit \texttt{0e1f7da},
\texttt{data/splits/}): \emph{dev} 80 tests (13/39/28 across Commons-cli/Gson/Jsoup), \emph{test} 150 (24/74/52) and \emph{rest} 254 (40/125/89).
Prompts, thresholds, attempt and token caps were tuned on dev only. All reported LLM runs use the test split and the runner's \texttt{--frozen}
guard, which refuses to start from a dirty tree and records the git commit. The rest split is used only for runs without an LLM (E1, E4 detector agreement).

\textbf{Systems.} E1 is UTRefactor's published outputs measured with our verifier (\textsf{b0replay}, no LLM).
E2 compares: \textsf{b0}, UTRefactor's prompt and knowledge base with our model, one pass; \textsf{b1}, the same prompt followed by a fixed retry
script (at most 2 retries with the error text); \textsf{single}, a plain prompt without the knowledge base; \textsf{b2}, a single tool-using agent
with the same tools and token cap as our system; and \textsf{s3}, our multi-agent system. All systems share the model, knowledge base, tools,
per-test token cap and test set. E3 is \textsf{s3} with all three gap components, E5 switches one off at a time
(\textsf{-G1}: kill-set measured, never rejects; \textsf{-G2}: tsDetect only; \textsf{-G6}: fixed retry policy; and memory off).
E2 numbers are labelled \emph{reproduced} when we ran the system and \emph{paper-reported} when they are copied from~\cite{utrefactor}.

\textbf{Model and parameters.} All LLM calls go to Cerebras \texttt{gpt-oss-120b} (free tier, 5 requests/min, 950k tokens/day) at temperature 0.7, at most
8000 output tokens per call and 60\,000 tokens per test. Every configuration on the test split is run with seeds 1, 2 and 3 (see \texttt{experiments/E*.json}), and we report mean $\pm$ sd over them; the replay \textsf{b0replay} calls no LLM and is run once. A sensitivity run (ES) varies the retry cap (0, 1, 2, 4) for \textsf{s3} on the first 20 dev tests; the LLM response cache is keyed by prompt and seed.
Every run records the model ID, temperature, system-prompt hash, seed and git commit in \texttt{config.json} and in the trajectory metadata.

\textbf{Tooling.} tsDetect (v3 jar from the replication package) and JNose (jnose-core 0.9.0, the last release built for Java~21) detect smells at method level;
PIT 1.17.4 computes the full mutation matrix; the fast runner compiles once per Maven module and runs each test class through the JUnit launcher.
The reported runs used an HP Omen~16 laptop with Ubuntu 24.04 under WSL2 (32 logical CPUs and 15~GB of memory visible to Linux), Python~3.12, OpenJDK~17 for the project builds, JDK~21 for JNose and Maven~3.9.9.
The pipeline restores every touched file after each test, and a post-test \texttt{git status} check stops the run if a checkout was left modified.

\textbf{Execution of the reported runs.} All 43 reported runs used the \texttt{--frozen} guard, which refuses to start when a tracked file differs from the commit and records the commit in each run's \texttt{config.json}: the 19 seed-1 runs (and seeds 2--3 of E3 and \textsf{-G1}) were made from commit \texttt{3b992f2}, the 24 seed-2 and seed-3 runs of E2, E5 and ES from \texttt{e9ff1c8}. The two commits share the prompts (same prompt hash), agents, verifier and metrics; the later one adds atomic cache writes, a separate label for runner timeouts, the \texttt{--input} option and the post-hoc JaCoCo check.
To fit the free tier, configurations ran in parallel, each in its own copy of the repository with its own API key, so no two runs ever shared a key or a working directory. A run that hit a key's daily limit was resumed with another key; resumption skips the tests whose results are already saved and redoes only the test in progress.
With up to twelve runs on one machine, PIT used 4 to 8 threads per run. This affects wall-clock times (Table~\ref{tab:e6}); outcomes could change only through PIT's timeout rule, since a slower run can turn a surviving mutant into a timed-out, and therefore killed, one.

